\uuid{y8P7}
\exo7id{5130}
\titre{exo7 5130}
\auteur{rouget}
\organisation{exo7}
\datecreate{2010-06-30}
\isIndication{false}
\isCorrection{true}
\chapitre{Nombres complexes}
\sousChapitre{Forme cartésienne, forme polaire}
\cpge{mpsi-mp2i.nombres-complexes.d, pcsi.nombres-complexes.d, ptsi.nombres-complexes.d, mpsi-mp2i.nombres-complexes.c, pcsi.nombres-complexes.c, ptsi.nombres-complexes.c, bcpst1.nombres-complexes.b}
\module{Algèbre}
\niveau{L1}
\difficulte{}

\contenu{
\question{
Calculer $(1+i\sqrt{3})^9$.
}
\reponse{
$(1+i\sqrt{3})^9=(2e^{i\pi/3})^9=2^9e^{3i\pi}=-512$.

\begin{center}
\fbox{
\begin{tabular}{l}
La forme algébrique d'un complexe est particulièrement bien adaptée à l'addition.\\
La forme trigonométrique d'un complexe est particulièrement bien adaptée à la multiplication.
\end{tabular}
}
\end{center}
}
}
